\documentclass[11pt,a4paper]{article}
\usepackage[utf8]{inputenc}
\usepackage[ngerman]{babel}
\usepackage{geometry}
\usepackage{amsmath}
\usepackage{outlines}
\geometry{a4paper, margin=2cm}
\usepackage{enumitem}

\title{Exposé: Aufgabenverteilung im WASP-Projekt}
\author{Ivan Kurilin \and Dimitrij Pivovar}
\date{\today}

\begin{document}	
	\maketitle		
	\section*{Aufgabenverteilung}
	
\subsection*{Ivan Kurilin}
\begin{itemize}
	\item \textbf{Stock mit Ballkollision:} Interaktion zwischen dem Stock und dem Ball, wie sie in der Simulation behandelt wird.
	\item \textbf{Kollision mit Banden:} Kollisionseffekte, die zwischen dem Ball und den Banden auftreten.
	\item \textbf{Billiard Aufstellung:} Objekte auf dem Billiardtisch zu Beginn des Spiels, um die Simulation zu starten.
\end{itemize}
	
\subsection*{Dimitrij Pivovar}
\begin{itemize}
	\item \textbf{3D Scene:} Die Szene aus verschiedenen Blickwinkeln betrachtbar.
	\item \textbf{2D Tisch-Erweiterung (Billiard Taschen, Banden):} Erweiterung des 2D-Tisches, die Billiardtaschen und Banden integriert, um eine realistische Interaktion der Bälle zu ermöglichen.
\end{itemize}

\subsection*{Mögliche Erweiterung (wird im Team aufgeteilt)}
\begin{itemize}
	\item \textbf{3D Bälle:} Modellierung der 3D-Bälle, einschließlich ihrer Bewegungen und Interaktionen auf dem Tisch.
	\item \textbf{Kugeldesign:} Das Design der Bälle, einschließlich Farbe, Zahlen, die im Billiard verwendet werden.
	\item \textbf{Regelsystem \& Punktensystem:} Die Definition der Regeln, die das Spiel steuern, einschließlich der Bedingungen für das Gewinnen und Verlieren.
\end{itemize}

	\section*{Zeitaufwand}
\begin{itemize}
	\item Geschätze Deadline: 10.03.2025.
	\item Geschätzer Zeitaufwand: 1-2 Wochen pro Thema.
\end{itemize}

\section*{Erwartete Ergebnisse}
\begin{itemize}
	\item Ein funktionierender 3D-Billard-Demonstrator.
	\item Unterstützung von Kollisionserkennung für verschiedene Objekte.
	\item Ein dokumentierter Bericht über die Umsetzung und die Herausforderungen.
\end{itemize}
	
\section{3D Scene }
\begin{itemize}
	\item Besteht aus einer Box, die als Tischoberfläche dient mit den Maßen x=300, y=30, z=200
	\item Beine mit je Maßen x=50, y=100, z=40. Mit 2 Schleifen implementiert:
	\begin{outline}
	\1 Äußere Schleife beginnt bei den Koordinaten x = -90, Kondition ist das x <= 90 ist und das Inkrement ist 180. Das sorgt dafür das die innere Schleife auf die linken (vom Startpunkt) Beine angewandt wird bevor die x Koordinate um 180 erhöht wird und dann die rechten Beine gezeichnet wird
	\1 Innere Schleife beginnt bei den Koordinaten x = -90, z = -40. z wird nach dem Schleifenkörper um 80 erhöht um auf die Position des hinteren Beins zu wechseln.
	\1 Zusammenfassend kann man sagen die Beine haben folgende Koordinaten:
	\2 Bein links hinten x = width/2 - 90, y = height/2+55, z = -40
	\2 Bein links vorne x = width/2 - 90, y = height/2+55, z = 40
	\2 Bein rechts hinten x = width/2 + 90, y = height/2+55, z = -40
	\2 Bein rechts vorne x = width/2 + 90, y = height/2+55, z = 40
	\end{outline}
	\item TODOs
	\begin{outline}
		\1 Bounding Box für Bälle auf Tischoberfläche begrenzen
		\1 Kameraposition verbessern für Startszene
		\1 Bälle in der Pyramide aufstellen für Startszene
		\1 Taschen implementieren
		\1 Kollision zwischen der BilliardCue und Bällen verbessern
		\1 Quellen + Formeln richtig Zitieren
	\end{outline}
	\item Herausforderungen
	\begin{outline}
		\1 Allgemeines Verständnis der translate() Funktion
		\2 Erstes Learning war, dass die translate Funktion akkumuliert. translate(50,0,0) und danach translate(70,0,0) summiert sich auf. pushMatrix(), popMatrix() hilft
	\end{outline}
	\end{itemize}
	\newpage
	\section{Kollisionsdetektion und -antwort}

Die Kollisionsdetektion zwischen den Bällen erfolgt, indem der Abstand zwischen den Mittelpunkten der beiden Bälle berechnet wird. Wenn dieser Abstand kleiner ist als die Summe der Radien der Bälle, wird eine Kollision erkannt. Die Berechnung des Abstands erfolgt mit der Formel:

\[
\text{Abstand} = \sqrt{(dx)^2 + (dy)^2}
\]

Dabei sind \(dx\) und \(dy\) die Differenzen der \(x\)- und \(y\)-Koordinaten der beiden Bälle.

\subsection{Zentraler elastischer Stoß}


Einen Stoß kann man als elastisch bezeichnen, falls sich die Energie der Objekte nach dem Stoß erhalten bleibt. Wir können den elastischen Stoß durch folgende Gleichung virtuell darstellen:

\[
m_1 v_1 + m_2 v_2 = m_1 v_1' + m_2 v_2'
\]

\begin{outline}
	\1 $m_1, m_2$: die Massen der beiden Objekte
	\1 $v_1, v_2$: die Geschwindigkeiten der beiden Objekte
\end{outline}

Die rechte Seite der Gleichung steht für den Impuls nach dem Stoß.



\subsection{Normalvektor und Impulsberechnung}
Der Normalvektor \((nx, ny)\) definiert die Richtung der Kollision. Er wird durch die Berechnung des Abstands zwischen den Bällen bestimmt.


\subsection{Placeholder}
	%Der Normalvektor und die Impulsberechnung sorgen für die korrekte Anwendung der Kollisionskräfte in der richtigen Richtung.


\end{document}
